\begin{figure}[t]
\centering
\begin{tikzpicture}[x=1cm,y=1cm]
\begin{scope}
\draw[thin] (0,2.4)--(0,0)--(6,0);
\foreach \x/\lab in {1.57143/1,2.82143/2,4.78571/3}{
 \draw[gray!50,thin] (\x,0)--(\x,.1862);
 \node[below] at (\x,-.02){$s_{\lab}$};}
\foreach \y in {.2,1.1,2.0}{\draw[gray!50,thin] (0,\y)--(.1766,\y);}
% Original printed PDF vector path.
\path[draw=plotgray,line width=.5pt] (0.50002,0.20002) .. controls (0.50002,0.20002) and (0.57375,0.20002) ..(0.60206,0.20002) .. controls (0.63038,0.20002) and (0.67579,0.20002) ..(0.70410,0.20002) .. controls (0.73242,0.20002) and (0.77783,0.20002) ..(0.80614,0.20002) .. controls (0.83446,0.20002) and (0.87987,0.20002) ..(0.90818,0.20002) .. controls (0.93650,0.20002) and (0.98191,0.19999) ..(1.01022,0.20002) .. controls (1.03854,0.20004) and (1.08395,0.19915) ..(1.11226,0.20022) .. controls (1.14058,0.20129) and (1.18599,0.19148) ..(1.21430,0.20774) .. controls (1.24262,0.22400) and (1.28803,0.21850) ..(1.31634,0.31741) .. controls (1.34466,0.41631) and (1.39007,0.68912) ..(1.41838,0.92059) .. controls (1.44670,1.15205) and (1.49211,1.83776) ..(1.52042,1.98565) .. controls (1.54874,2.13353) and (1.59415,2.13408) ..(1.62246,1.98645) .. controls (1.65078,1.83881) and (1.69619,1.15313) ..(1.72450,0.92159) .. controls (1.75282,0.69005) and (1.79823,0.41673) ..(1.82654,0.31769) .. controls (1.85486,0.21865) and (1.90027,0.22406) ..(1.92858,0.20776) .. controls (1.95690,0.19147) and (2.00231,0.20130) ..(2.03062,0.20022) .. controls (2.05894,0.19915) and (2.10434,0.20004) ..(2.13266,0.20002) .. controls (2.16098,0.19999) and (2.20639,0.20001) ..(2.23470,0.20002) .. controls (2.26302,0.20002) and (2.30842,0.19962) ..(2.33674,0.20009) .. controls (2.36506,0.20056) and (2.41046,0.19445) ..(2.43878,0.20340) .. controls (2.46710,0.21236) and (2.51250,0.19603) ..(2.54082,0.26464) .. controls (2.56914,0.33326) and (2.61454,0.49200) ..(2.64286,0.69792) .. controls (2.67118,0.90385) and (2.71658,1.54801) ..(2.74490,1.74880) .. controls (2.77322,1.94960) and (2.81862,2.22316) ..(2.84694,2.14511) .. controls (2.87525,2.06706) and (2.92066,1.42815) ..(2.94898,1.18627) .. controls (2.97730,0.94440) and (3.02270,0.53638) ..(3.05102,0.40186) .. controls (3.07933,0.26733) and (3.12474,0.24463) ..(3.15306,0.21670) .. controls (3.18137,0.18877) and (3.22678,0.20289) ..(3.25510,0.20057) .. controls (3.28341,0.19826) and (3.32882,0.20010) ..(3.35714,0.20002) .. controls (3.38545,0.19995) and (3.43086,0.20002) ..(3.45918,0.20002) .. controls (3.48749,0.20002) and (3.53290,0.20002) ..(3.56122,0.20002) .. controls (3.58953,0.20002) and (3.63494,0.20002) ..(3.66326,0.20002) .. controls (3.69157,0.20002) and (3.73698,0.20002) ..(3.76530,0.20002) .. controls (3.79361,0.20002) and (3.83902,0.20002) ..(3.86734,0.20002) .. controls (3.89565,0.20002) and (3.94106,0.20002) ..(3.96938,0.20002) .. controls (3.99769,0.20002) and (4.04310,0.20002) ..(4.07142,0.20002) .. controls (4.09973,0.20002) and (4.14514,0.20001) ..(4.17346,0.20002) .. controls (4.20177,0.20002) and (4.24718,0.19984) ..(4.27550,0.20004) .. controls (4.30381,0.20023) and (4.34922,0.19676) ..(4.37754,0.20142) .. controls (4.40585,0.20607) and (4.45126,0.18872) ..(4.47958,0.23359) .. controls (4.50789,0.27847) and (4.55330,0.35351) ..(4.58162,0.52488) .. controls (4.60993,0.69625) and (4.65534,1.23631) ..(4.68366,1.46874) .. controls (4.71197,1.70117) and (4.75738,2.19939) ..(4.78570,2.20005) .. controls (4.81401,2.20072) and (4.85942,1.70563) ..(4.88773,1.47354) .. controls (4.91605,1.24145) and (4.96146,0.69929) ..(4.98978,0.52730) .. controls (5.01809,0.35531) and (5.06350,0.27920) ..(5.09182,0.23398) .. controls (5.12013,0.18877) and (5.16554,0.20615) ..(5.19386,0.20144) .. controls (5.22217,0.19673) and (5.26758,0.20024) ..(5.29589,0.20004) .. controls (5.32421,0.19984) and (5.36962,0.20002) ..(5.39793,0.20002) .. controls (5.42625,0.20001) and (5.49998,0.20002) ..(5.49998,0.20002);

\node[below] at (3,-.65){$x$};
\node[rotate=90] at (-.42,1.2){$p(x)$};
\end{scope}
\begin{scope}[xshift=7.4cm]
\draw[thin] (0,2.4)--(0,0)--(6,0);
\foreach \x/\lab in {1.57143/1,2.82143/2,4.78571/3}{
 \draw[gray!50,thin] (\x,0)--(\x,.1862);
 \node[below] at (\x,-.02){$s_{\lab}$};}
\foreach \y in {.2,1.1,2.0}{\draw[gray!50,thin] (0,\y)--(.1766,\y);}
% Original printed PDF vector path.
\path[draw=plotgray,line width=.5pt] (0.50002,0.20000) .. controls (0.50002,0.20000) and (0.57375,0.19999) ..(0.60206,0.20006) .. controls (0.63038,0.20014) and (0.67579,0.20008) ..(0.70410,0.20055) .. controls (0.73242,0.20101) and (0.77783,0.20116) ..(0.80614,0.20339) .. controls (0.83446,0.20562) and (0.87987,0.20809) ..(0.90818,0.21661) .. controls (0.93650,0.22512) and (0.98191,0.23915) ..(1.01022,0.26477) .. controls (1.03854,0.29038) and (1.08395,0.34105) ..(1.11226,0.40122) .. controls (1.14058,0.46140) and (1.18599,0.58978) ..(1.21430,0.69847) .. controls (1.24262,0.80715) and (1.28803,1.03871) ..(1.31634,1.18456) .. controls (1.34466,1.33042) and (1.39007,1.61653) ..(1.41838,1.74970) .. controls (1.44670,1.88286) and (1.49211,2.08957) ..(1.52042,2.14436) .. controls (1.54874,2.19914) and (1.59415,2.19925) ..(1.62246,2.14458) .. controls (1.65078,2.08990) and (1.69619,1.88339) ..(1.72450,1.75028) .. controls (1.75282,1.61717) and (1.79823,1.33101) ..(1.82654,1.18520) .. controls (1.85486,1.03940) and (1.90026,0.80787) ..(1.92858,0.69943) .. controls (1.95690,0.59098) and (2.00231,0.46234) ..(2.03062,0.40363) .. controls (2.05894,0.34491) and (2.10434,0.29567) ..(2.13266,0.27625) .. controls (2.16098,0.25684) and (2.20639,0.25232) ..(2.23470,0.26368) .. controls (2.26302,0.27504) and (2.30842,0.31062) ..(2.33674,0.35815) .. controls (2.36505,0.40567) and (2.41046,0.51046) ..(2.43878,0.60621) .. controls (2.46710,0.70195) and (2.51250,0.90853) ..(2.54082,1.04820) .. controls (2.56914,1.18788) and (2.61454,1.47022) ..(2.64286,1.61287) .. controls (2.67117,1.75552) and (2.71658,1.99677) ..(2.74490,2.07634) .. controls (2.77322,2.15591) and (2.81862,2.21416) ..(2.84694,2.18637) .. controls (2.87525,2.15859) and (2.92066,1.99527) ..(2.94898,1.87609) .. controls (2.97730,1.75691) and (3.02270,1.47610) ..(3.05102,1.32742) .. controls (3.07933,1.17873) and (3.12474,0.92508) ..(3.15306,0.80448) .. controls (3.18137,0.68388) and (3.22678,0.52988) ..(3.25510,0.45822) .. controls (3.28341,0.38656) and (3.32882,0.32052) ..(3.35714,0.28801) .. controls (3.38545,0.25550) and (3.43086,0.23539) ..(3.45918,0.22390) .. controls (3.48749,0.21240) and (3.53290,0.20836) ..(3.56122,0.20517) .. controls (3.58953,0.20198) and (3.63494,0.20159) ..(3.66326,0.20089) .. controls (3.69157,0.20019) and (3.73698,0.20023) ..(3.76530,0.20014) .. controls (3.79361,0.20004) and (3.83902,0.20004) ..(3.86734,0.20021) .. controls (3.89565,0.20039) and (3.94106,0.20036) ..(3.96938,0.20140) .. controls (3.99769,0.20244) and (4.04310,0.20323) ..(4.07142,0.20770) .. controls (4.09973,0.21218) and (4.14514,0.21847) ..(4.17346,0.23367) .. controls (4.20177,0.24887) and (4.24718,0.27676) ..(4.27550,0.31724) .. controls (4.30381,0.35773) and (4.34922,0.44181) ..(4.37754,0.52544) .. controls (4.40585,0.60906) and (4.45126,0.78890) ..(4.47958,0.91994) .. controls (4.50789,1.05097) and (4.55330,1.32204) ..(4.58162,1.46983) .. controls (4.60993,1.61762) and (4.65534,1.88383) ..(4.68366,1.98512) .. controls (4.71197,2.08642) and (4.75738,2.19967) ..(4.78569,2.19989) .. controls (4.81401,2.20011) and (4.85942,2.08766) ..(4.88774,1.98670) .. controls (4.91605,1.88574) and (4.96146,1.61994) ..(4.98978,1.47223) .. controls (5.01809,1.32452) and (5.06350,1.05332) ..(5.09181,0.92212) .. controls (5.12013,0.79092) and (5.16554,0.61048) ..(5.19386,0.52664) .. controls (5.22217,0.44279) and (5.26758,0.35843) ..(5.29589,0.31781) .. controls (5.32421,0.27719) and (5.36962,0.24913) ..(5.39793,0.23386) .. controls (5.42625,0.21860) and (5.49997,0.20775) ..(5.49997,0.20775);

\node[below] at (3,-.65){$y$};
\node[rotate=90] at (-.42,1.2){$p(y)$};
\end{scope}
\end{tikzpicture}
\caption{一个由 $X$ 导致 $Y$ 的例子。由此，$P_Y$ 与 $P_{X\mid Y}$ 包含关于彼此的信息。左：$P_X$ 是由位于 $s_1,s_2,s_3$ 的尖峰组成的混合分布。右：$P_Y$ 由 $P_X$ 与零均值高斯噪声进行卷积（convolution）得到，因此，其峰仍位于 $s_1,s_2,s_3$，但不再那么尖锐。于是，$P_{X\mid Y}$ 同样包含关于 $s_1,s_2,s_3$ 的信息，见习题 5.1。}
\label{fig:ch05-shared-peak-information}
\end{figure}
